 \vspace{0.3cm}
  \begin{resumo}
     %%%%%%%%%%%%%% Não retirar este trecho
    \noindent\footnote{Trabalho de Conclusão de Curso apresentado ao Curso de Engenharia de Computação no segundo semestre de 2025, Área de Ciências Exatas e Tecnológicas da Universidade do Oeste de Santa Catarina como requisito parcial à obtenção do grau de bacharel em Engenharia de Computação.}
     %%%%%%%%%%%%%% Não retirar este trecho
     Este trabalho propõe o desenvolvimento de uma interface gráfica para telemetria e parametrização de robôs autônomos, integrada ao \textit{Robot Operating System} (ROS) por meio do protocolo ROSBridge. Utilizando a biblioteca Streamlit, a interface foi projetada para oferecer interação intuitiva, com visualização em tempo real de velocidade, posição e mapas 2D gerados por algoritmos SLAM (\textit{Simultaneous Localization and Mapping}), além de possibilitar ajustes de parâmetros operacionais e controle manual do robô. A metodologia empregada estrutura-se em três fases distintas: a caracterização do \textit{hardware}, o desenvolvimento modular e a validação em ambientes simulados (RViz) e reais. O desenvolvimento modular encontra-se concluído, com o módulo de visualização e de controle implementados. \textit{Scripts} de comunicação via ROSBridge estão operacionais e integrados à interface e funcionalidades como o mapeamento 2D com LiDAR (Light Detection and Ranging) e a implementação de \textit{waypoints} estão funcionando. Funcionalidades planejadas na primeira fase de desenvolvimento do trabalho de conclusão de curso, como plotar mapa 2D com nuvem de pontos e navegação autônoma por waypoints foram implementadas com sucesso, testes também foram realizados para validar a precisão de posicionamento no mapa e a eficácia da interface. O sistema demonstra potencial para aplicações em ambientes internos e semi-estruturados, contribuindo para a integração eficiente entre interfaces gráficas e sistemas robóticos.
     
  \end{resumo}
  
  \begin{palavraschave}
    aplicativo, app, interface gráfica, LiDAR, robô autônomo, ROSBridge, Streamlit, web.
  \end{palavraschave}
  
  \vspace{0.5cm}
  
  \begin{abstract}
  \noindent This work proposes the development of a graphical interface for telemetry and parameterization of autonomous robots, integrated with the \textit{Robot Operating System} (ROS) through the ROSBridge protocol. Using the Streamlit library, the interface was designed to offer intuitive interaction, with real-time visualization of velocity, position, and 2D maps generated by SLAM (\textit{Simultaneous Localization and Mapping}) algorithms, in addition to enabling adjustments of operational parameters and manual control of the robot. The methodology employed is structured in three distinct phases: \textit{hardware} characterization, modular development, and validation in simulated (RViz) and real environments. The modular development is concluded, with the visualization and control modules implemented. \textit{Scripts} for communication via ROSBridge are operational and integrated into the interface, and functionalities such as 2D mapping with LiDAR and the implementation of \textit{waypoints} are working. Planned features for the first development phase of the undergraduate thesis, such as plotting 2D maps with point clouds and autonomous navigation by waypoints were implemented successfully, tests were also conducted to validate the positioning accuracy on the map and the interface's effectiveness. The system demonstrates potential for applications in indoor and semi-structured environments, contributing to the efficient integration between graphical interfaces and robotic systems.
  \end{abstract}

  \begin{IEEEkeywords}
    application, app, Autonomous robot, graphical interface, LiDAR, ROSBridge, Streamlit, web.
  \end{IEEEkeywords}
  
    \vspace{0.5cm}